% ============================================================
% ch01_introduction.tex
% Chapters: Introduction, Problem Statement, Objectives,
%           Literature Review, Existing vs Proposed System
% ============================================================

\chapter{Introduction}

\section{Background}

Modern database management systems must handle unexpected failures gracefully.
When a system crashes mid-transaction, data can be left in an inconsistent
state --- some changes applied, others lost. The mechanism that prevents
permanent data loss is called \textbf{database recovery}, and it relies on
\textbf{transaction logs} to determine what was committed and what was not.

This project implements a simplified \textbf{Log-Based Recovery System} that
demonstrates these concepts using two databases working together:

\begin{itemize}
  \item \textbf{PostgreSQL} --- a relational database storing structured
        banking data (accounts, customers, transactions).
  \item \textbf{MongoDB} --- a document database storing Write-Ahead Log
        (WAL) entries that record every transaction before it is applied.
\end{itemize}

The combination of two database technologies in one system is called
\textbf{polyglot persistence} --- a core concept in modern ADBMS.

\section{Motivation}

Traditional textbook explanations of database recovery are often abstract.
This project makes recovery concepts tangible by:

\begin{itemize}
  \item Simulating real crashes during active transactions
  \item Showing exactly what an UNCOMMITTED log looks like in MongoDB
  \item Demonstrating how \textit{old\_balance} and \textit{new\_balance}
        in a WAL log enable Undo and Redo operations
  \item Running a recovery process that restores database consistency
\end{itemize}

The banking domain (deposits, withdrawals, transfers) provides a familiar
scenario where inconsistency --- such as money being deducted without being
credited --- is easy to understand and verify.

\section{Scope of the Project}

\begin{conceptbox}
\textbf{This project is educational, not production-grade.}
The goal is to demonstrate ADBMS concepts clearly at undergraduate level,
not to implement a full DBMS recovery engine.
\end{conceptbox}

\noindent The scope includes:
\begin{itemize}
  \item Six PostgreSQL tables for banking data
  \item Three MongoDB collections for WAL and recovery logs
  \item Python modules for transactions, crash simulation, and recovery
  \item A simplified three-phase recovery process inspired by ARIES
\end{itemize}

The scope excludes: full ARIES internals, LSN chains, buffer pool management,
distributed recovery, and multi-user concurrency control.

% ────────────────────────────────────────────────────────────
\chapter{Problem Statement}

Traditional relational databases are vulnerable during system crashes.
When a database process is interrupted mid-transaction, two types of
problems arise:

\begin{enumerate}
  \item \textbf{Dirty writes} --- changes applied to disk that were never
        committed (should be undone).
  \item \textbf{Lost commits} --- changes that were committed logically but
        not yet flushed to disk before the crash (should be redone).
\end{enumerate}

Without a recovery mechanism, the database is left in an inconsistent state.
For a banking system, this could mean:

\begin{itemize}
  \item A withdrawal deducted from the sender's account but the receiver's
        credit was never applied.
  \item A deposit appearing committed in the application but lost from
        the database storage.
  \item Balance values that match neither the pre-transaction state nor
        the post-transaction state.
\end{itemize}

\noindent\textbf{The central problem this project solves:}

\begin{conceptbox}
How can a database system detect and repair inconsistencies caused by a
crash, using only the information available in transaction log records?
\end{conceptbox}

% ────────────────────────────────────────────────────────────
\chapter{Objectives}

The primary objectives of this project are:

\begin{enumerate}
  \item \textbf{Implement Write-Ahead Logging (WAL)} --- Every transaction
        writes a log record to MongoDB \textit{before} applying changes to
        PostgreSQL.

  \item \textbf{Simulate Crashes} --- Inject a system crash mid-transaction
        to create an inconsistent database state for recovery testing.

  \item \textbf{Implement UNDO Recovery} --- Identify UNCOMMITTED transactions
        from WAL logs and restore \textit{old\_balance} in PostgreSQL.

  \item \textbf{Implement REDO Recovery} --- Identify COMMITTED transactions
        whose PostgreSQL changes were lost and re-apply \textit{new\_balance}.

  \item \textbf{Demonstrate Polyglot Persistence} --- Use PostgreSQL and
        MongoDB together, each for what it does best.

  \item \textbf{Provide a Practical Demo} --- Create a menu-driven Python
        application that demonstrates all concepts interactively.
\end{enumerate}

\noindent\textbf{ADBMS concepts covered:}

\begin{table}[H]
\centering
\caption{ADBMS Concepts and Their Implementation}
\begin{tabularx}{\textwidth}{lX}
\toprule
\textbf{Concept} & \textbf{How Demonstrated} \\
\midrule
Transactions & Python TransactionManager with explicit COMMIT/ROLLBACK \\
ACID Properties & Atomicity via rollback; Durability via WAL \\
Write-Ahead Logging & MongoDB log written before every PostgreSQL update \\
Undo Recovery & old\_balance restored for UNCOMMITTED logs \\
Redo Recovery & new\_balance re-applied for lost COMMITTED changes \\
Crash Recovery & CrashSimulator raises exception before COMMIT \\
Polyglot Persistence & PostgreSQL for data + MongoDB for logs \\
\bottomrule
\end{tabularx}
\end{table}

% ────────────────────────────────────────────────────────────
\chapter{Literature Review}

\section{Database Transaction Management}

A \textbf{transaction} is a logical unit of work that must be executed
completely or not at all. The ACID properties --- Atomicity, Consistency,
Isolation, and Durability --- define the guarantees that a reliable
transaction system must provide \cite{ramakrishnan2002}.

\begin{itemize}
  \item \textbf{Atomicity}: A transaction either completes fully or has
        no effect. PostgreSQL achieves this through \texttt{COMMIT} and
        \texttt{ROLLBACK}.
  \item \textbf{Consistency}: The database moves from one valid state to
        another. CHECK constraints and foreign keys enforce this.
  \item \textbf{Isolation}: Concurrent transactions do not interfere.
        PostgreSQL uses MVCC (Multi-Version Concurrency Control).
  \item \textbf{Durability}: Committed changes survive crashes. WAL logging
        ensures durability by writing changes to a log before applying them.
\end{itemize}

\section{Write-Ahead Logging}

Write-Ahead Logging (WAL) is the fundamental technique for ensuring database
durability \cite{mohan1992}. The WAL rule states:

\begin{conceptbox}
\textbf{WAL Rule}: A log record describing a data change must be written to
stable storage \textit{before} the data change itself is written.
\end{conceptbox}

In our system, ``stable storage'' is MongoDB and ``data change'' is the
PostgreSQL balance update. Every transaction follows this sequence:
\begin{enumerate}
  \item Write log to MongoDB (UNCOMMITTED)
  \item Update PostgreSQL
  \item COMMIT PostgreSQL
  \item Update MongoDB log to COMMITTED
\end{enumerate}

\section{ARIES Recovery Algorithm}

The ARIES (Algorithm for Recovery and Isolation Exploiting Semantics) algorithm
\cite{mohan1992} is the industry standard for crash recovery in relational
databases. It operates in three phases:

\begin{enumerate}
  \item \textbf{Analysis Phase}: Scan the log to identify committed and
        uncommitted transactions at crash time.
  \item \textbf{Redo Phase}: Re-apply all logged changes to ensure the
        database reflects all committed operations.
  \item \textbf{Undo Phase}: Reverse all changes made by uncommitted
        transactions.
\end{enumerate}

Our project implements a simplified version of this three-phase approach
without LSN chains, dirty page tables, or checkpoint records.

\section{Polyglot Persistence}

Polyglot persistence refers to using multiple database technologies within
a single application, each chosen for its strengths \cite{fowler2012}.
NoSQL databases like MongoDB are well-suited for flexible, high-speed log
storage because:

\begin{itemize}
  \item Schema-free documents can store variable metadata per log entry
  \item Insert operations are fast and do not require transaction overhead
  \item Rich query capabilities allow efficient log filtering by status
  \item MongoDB is independent of the PostgreSQL transaction — a failed
        PG transaction does not affect the MongoDB log
\end{itemize}

% ────────────────────────────────────────────────────────────
\chapter{Existing System vs Proposed System}

\section{Existing Approaches}

Traditional database recovery approaches have significant limitations for
educational purposes:

\begin{table}[H]
\centering
\caption{Existing vs Proposed System Comparison}
\begin{tabularx}{\textwidth}{lXX}
\toprule
\textbf{Aspect} & \textbf{Existing/Traditional} & \textbf{Proposed System} \\
\midrule
Log Storage & Binary log files on disk & MongoDB JSON documents \\
Log Format & Internal binary format & Human-readable JSON \\
Recovery & Automated, invisible to user & Explicit, demonstrable steps \\
Complexity & Full ARIES with LSNs & Simplified 3-phase approach \\
Observability & Cannot inspect logs easily & Query logs with mongosh \\
Database Type & Single RDBMS & Hybrid PostgreSQL + MongoDB \\
Educational Use & Not designed for teaching & Specifically for ADBMS demos \\
Crash Testing & Requires actual system crash & Python Exception injection \\
\bottomrule
\end{tabularx}
\end{table}

\section{Advantages of the Proposed System}

\begin{enumerate}
  \item \textbf{Transparency}: Every log entry is a readable JSON document
        that students can inspect and understand.
  \item \textbf{Demonstrability}: The crash and recovery sequence is
        observable step-by-step in the terminal.
  \item \textbf{Concept Clarity}: The separation of WAL (MongoDB) and data
        (PostgreSQL) makes the WAL principle visually clear.
  \item \textbf{Practical Implementation}: Runs on a single laptop with
        freely available software.
\end{enumerate}
